% !TeX root = ../main.tex

\chapter{文獻回顧}
\label{chap:related}

本章回顧與本論文研究直接相關之四類文獻：第~\ref{sec:ienf-related}~節說明 IENFD 量化之臨床背景與既有自動化嘗試；第~\ref{sec:vesselness}~節整理影像增強與血管／纖維性濾波之理論基礎；第~\ref{sec:segmentation}~節討論既有 IENF 切分方法之兩條路徑，及其與拓樸層計數之間之轉換缺口；第~\ref{sec:graph-tracing}~節回顧圖論式纖維追蹤，特別是最短路徑、最小生成樹與骨架化方法在生醫影像重建中的應用。

\section{表皮內神經纖維密度量化}
\label{sec:ienf-related}

IENFD 自 2005 年 EFNS 發佈之量測指引（後於 2010 年更新）~\cite{lauria2010skinbiopsy, jointtaskforce2010skinbiopsy} 以來被奉為小纖維神經病變診斷之參考標準，其量測規則要求判讀者於 PGP 9.5 染色之皮膚垂直切片中，沿表皮／真皮交界逐一辨識由真皮穿越進入表皮之神經纖維，並將「一條進入後分支者」僅計為一次。由於此規則同時涉及「拓樸連續性判定」與「跨越事件去重」兩項子任務，早期的自動化嘗試多侷限在第一步之影像增強與分割，並未真正觸及拓樸層面之推理。

部份研究~\cite{vanacker2016automated} 提出以形態學閉運算配合 Otsu 閾值對 PGP 9.5 影像進行二值化，再以連通元件數粗略估計纖維數，但由於缺乏跨越去重機制，容易高估密度，亦無法處理表皮內分支。類似策略亦見於角膜共軛焦顯微 (CCM) 影像之糖尿病周邊神經病變評估~\cite{dabbah2011automatic}。另一類方法則以半自動化方式讓使用者先標出表皮邊界，再以距離場與局部峰值偵測計算跨越點，雖能得到可重複之結果，但仍需要人工輸入關鍵資訊。上述研究皆顯示，從分割結果到可計數之拓樸結構之間存在結構性缺口，而此缺口正是本研究所欲處理之核心。

\section{影像增強與血管／纖維性濾波}
\label{sec:vesselness}

PGP 9.5 染色呈現之神經纖維在影像上為細長之線狀結構 (curvilinear structure)，其灰階變化在局部呈現「一個方向強變化、另一方向弱變化」的二階微分特徵。Frangi 等人 (1998)~\cite{frangi1998multiscale} 將此觀察形式化為基於 Hessian 矩陣特徵值之血管性 (vesselness) 量測：設 $\lambda_1, \lambda_2$ 為影像於多尺度 Gaussian 平滑後之 Hessian 特徵值且 $|\lambda_1| \le |\lambda_2|$，則沿血管方向 $\lambda_1 \approx 0$、橫向 $|\lambda_2|$ 大，可據此設計出對線狀結構具高響應、對團塊或雜訊具低響應之濾波。

Sato 等人 (1998)~\cite{sato1998three} 進一步提出以不同權重加總 $\lambda_2$ 並對黑／白線同時處理之變體，其輸出對多尺度結構（粗細不一之纖維）具有良好之穩定性，並廣泛實作於 \texttt{scikit-image}~\cite{vanderwalt2014scikit} 之 \texttt{sato} 函式中。

在血管性濾波之前，另需處理 PGP 9.5 染色常見之背景不均。研究上常以形態學開運算 (morphological opening)~\cite{soille1999morphological} 估計影像之慢變背景，再以原影像減去背景得到對比一致之結果；之後再以對比受限直方圖均衡化 (CLAHE)~\cite{zuiderveld1994contrast} 進一步提升局部對比。相較於對整張影像共用同一灰階重映射之全域直方圖均衡化，CLAHE 透過分塊適應與每塊裁剪閾值機制，可在強化細小纖維對比之同時抑制近均勻區域中之雜訊過度放大。本論文之前處理鏈即沿用此策略。

\section{神經纖維之語意分割與其破碎化問題}
\label{sec:segmentation}

於自動化 IENFD 量化之文獻中，現有切分方法可依其與拓樸資訊之關係粗分為兩條路徑，差異在於拓樸是於切分中內化處理，抑或延後至切分之後處理。

第一條路徑透過深度學習從原始影像端到端切分出神經纖維結構，已應用於 PGP 9.5 皮膚切片~\cite{chan2022automated} 與角膜共軛焦顯微影像~\cite{zhang2020automatic}，其輸出旨在直接還原一條條完整之纖維。此類方法將拓樸資訊內化於模型之中：是否跨越染色空隙、是否屬於同一條纖維等判定皆隱含於切分結果，而非以顯式步驟處理。其代價是訓練監督之門檻較高——通常需於像素層提供已解析至「纖維」層級之金標準，而非僅標出染色可見區段；切分失誤亦難以拆解為像素響應不足或拓樸連接錯誤，使各階段之優化難以獨立進行。

第二條路徑則聚焦於更基礎之 \textbf{IENF profile 切分}~\cite{huang2026automated}，即僅切分出 PGP 9.5 染色於切片上實際呈現之\textbf{纖維微粒} (particle)~\cite{hsueh2026machine}，而不嘗試於切分階段還原其拓樸。此條路徑以染色之破碎本質為前提：IENF profile 本即由多段不連續之纖維微粒組成，並不直接對應「一條神經纖維」之拓樸實體；同一條神經穿越表皮／真皮交界時，常以多個纖維微粒間斷地呈現於影像上。然而 EFNS 規範之有效跨越計數要求判定「一條進入後分支者僅計為一次」，此乃拓樸層之資訊，無法僅由像素層之 profile 切分得到。

本研究即聚焦於從切分到拓樸之轉換缺口，將 IENF profile 之切分結果視為可信但不完整之標註，並以圖論方法在保留可信微粒之前提下，補足微粒間缺失之拓樸連接，將像素級切分結果整合為可直接套用臨床計數規則之神經纖維結構。此設計與第一條路徑（端到端 DL 直接輸出纖維）具本質上的差異：本研究將「切分」與「拓樸重建」明確分離，使每一階段之失敗模式與設計取捨皆可獨立分析與優化。

\section{圖論方法於纖維重建}
\label{sec:graph-tracing}

圖論式纖維追蹤可追溯至血管與神經影像之最短路徑分析。Cohen 與 Kimmel 於 1997 年~\cite{cohen1997global} 提出以最小作用量原則 (minimal action principle) 從影像成本圖抽取最短路徑之方法，此後該框架廣被用於視網膜血管、腦神經纖維與顯微影像之重建。其核心想法是：將影像依據「是否屬於目標結構」轉換為「成本」，使沿目標行進成本低、橫跨背景成本高，則任兩點之最小成本路徑即為最可能之生物連接。

將該想法推廣至多個來源時，即形成\textbf{多源 Dijkstra 擴張}。若將每個連通元件視為一個來源並同步擴張，所得之 Voronoi-like 場~\cite{felzenszwalb2012distance} 可同時表達「每個像素最接近哪個元件」與「抵達該像素之最小成本」，並在兩元件之擴張前線相遇之位置找到其最佳連接。此種影像導引之元件連結概念與本文之 AEL 一脈相承，亦見於角膜共軛焦顯微影像之神經纖維量化~\cite{chen2016automatic}；但過往研究多針對血管重建、神經電生理追蹤或角膜纖維量化，並未針對「分段已有標註、需整合跨越判定」之 IENF 任務做系統性設計。

另一類方法則以\textbf{最小生成樹 (MST)} 連結既有微粒。Tagliasacchi 等人~\cite{tagliasacchi20163d} 在三維骨架抽取中展示 MST 能有效將點雲連接為無環之樹狀結構；在二維纖維影像上，亦有研究以微粒端點為節點、以歐式距離或局部方向差異為邊權建 MST。然而單獨使用 MST 會受到「距離短但生物上不合理」之邊支配，產生跨越真皮之錯接；若不搭配影像感知之邊權與適當之裁剪機制，MST 之結果往往需要大量後處理方能實用。

在獲得纖維重建之像素遮罩後，進一步之片段偵測與區域標記需操作於單像素寬之拓樸結構。Zhang 與 Suen~\cite{zhang1984fast} 提出之迭代細化演算法廣為此目的所用：其以鄰域連通性條件逐步剝除邊界像素，使輸入於拓樸不變之前提下被細化為中軸線，便於以節點度數抽取端點、分支點與片段。本論文以此演算法將整合後之纖維遮罩轉為最終之像素級拓樸圖。

本論文結合上述兩條技術脈絡：以多源 Dijkstra 於影像成本圖上所產生之相遇點作為候選邊，並以此建立元件層圖，再以成本閾值裁剪後求最小生成森林，同時利用 Dijkstra 所記錄之回溯指標將每條 MST 邊延伸為具體之像素路徑，最後與原標註合併並以 Zhang--Suen 細化為單像素寬之骨架，輸出可直接供下游分段與區域標記操作之像素級拓樸圖。此設計同時享有 Dijkstra 的影像感知性與 MST 的全域一致性。
